\chapter{Viertes Kapitel}
\section{Listen und Aufz\"ahlungen}

Hier mal eine Auflistung von Elementen
\begin{itemize}
 \item erstes Element
 \item zweites Element
 \item noch ein Element
\end{itemize}

Hier mal eine Aufz\"ahlung
\begin{enumerate}
 \item erster Punkt
 \item noch ein Punkt
 \item letzter Punkt
\end{enumerate}

\section{Und n\"achster Abschnitt etwas l\"anger als vorher es war}
Eine neue Seite, um auchmal die Kopfzeile zu sehen, da sie auf Seiten mit Kapitelanfang nicht erscheinen. Eine Abk\"urzung ist z.B. etc.\abk{etc.}{et cetera}.

\section{Eine Tabelle}

Hier eine Tabelle:
\begin{table}[htbp]
\centering
\begin{tabular}{l|l|l|l}
SpalteA & SpalteB & SpalteC & SpalteD \\
\midrule
InhaltA1 & InhaltB1 & InhaltC1 & InhaltD1 \\
InhaltA2 & InhaltB2 & InhaltC2 & InhaltD2 \\
InhaltA3 & InhaltB3 & InhaltC3 & InhaltD3
\end{tabular}
\caption{Beispiel einer Tabelle}
\label{tab:tabelle1}
\end{table}

Wie man in der Tabelle~\ref{tab:tabelle1} sehen kann ...

% Neue Seite
\section{Zitieren}
Zitierregeln sind der Word-Version der Diplomarbeitsvorlage zu entnehmen! \\
In der Datei \textit{bib.bib} im Verzeichnis \textit{Literatur} sind neue Eintr\"age von Literatur hinzuzuf\"ugen. Das Format kann in der \LaTeX{} Dokumentation nachgesehen werden. \\
N\"utzliche Links zum automatischen Erstellen von BibTeX Eintr\"agen:
\begin{itemize}
	\item http://truben.no/latex/bibtex/
	\item http://www.ottobib.com
	\item http://www.literatur-generator.de
	\item https://scholar.google.de
\end{itemize}
Damit das Literaturverzeichnis erstellt wird, muss auch mit dem Befehl BibTeX kompiliert werden! \\
\subsection{Beispiele}
Die exakte Herleitung kann in \cite{Prager1961} nachgelesen werden. \\
Hier ein Zitat etwas einger\"uckt:
\begin{quote}
Das Programm TeX wurde von Donald E. Knuth, Professor an der Stanford University, entwickelt. Leslie Lamport entwickelte Anfang der 1980er Jahre darauf aufbauend LaTeX, eine Sammlung von TeX-Makros. Der Name ist eine Abk\"urzung für Lamport TeX \cite{latextug}.
\end{quote}

\section{Bilder und Referenzen}
\begin{figure}[htbp]
\centering
\includegraphics[width=0.7\textwidth]{abbildungsdatei} % Datei in "bilder/" bei LaTeX: eps, bei PDFLaTeX: jpg (o.�.) 
\caption{Titel der Abbildung} 
\label{fig:bild1}
\end{figure}

In der Abbildung~\ref{fig:bild1}\footnote{vgl. Zitat A\cite{referenzA}} % referenzA muss im Verzeichnis literatur in der Datei bib.bib enthalten sein
ist zu sehen, dass ...

\section{Source Code einbinden}
Hier eine Einbindung von Source Code, funktioniert mit allen g\"angigen Programmiersprachen:
\lstinputlisting{./src/HelloWorld.java}
